% Transcription — fonds Alexandre Grothendieck, shelfmark 32, batch 1 (pages 1-20).
% Established from the facsimile archives/batches/32/batch-01.pdf (a mirror of
% the folder's PDF; no cover sheet, so PDF page k = folder page k), per the
% skill transcribe-grothendieck.
%
% Pass: Opus 5.5 (claude-opus-5-5), 2026-09-24 — first pass, unchecked against the pages by a human.
%
% The batch holds four pieces of correction material, three of them opened by
% a cover or a heading in his hand.
%
%   Page 2 — ink, top left of a blank leaf: « SGA 4 IX », a list of page and
%     line numbers to correct (pp. 2, 3, 60, 61, 63, 80, 91), no text.
%   Page 4 — pencil, under the orange cover of page 3 (« Notes et lettres »
%     struck, « XVII » boxed, « 5.5 »? in pencil): derived functors LT of a
%     non-additive functor T (Γ^n, Λ^n, Sym^n) via Dold–Puppe — commutation
%     with lim→, tor-dimension bounds, amplitude of N T N^{-1} K for K in
%     [p,q] and deg T ≤ n; ends « Ajouter remarque formulaire ».
%   Pages 6-9, 11-14 — typescript under the cover of page 5 (« Notes »,
%     « Bibliographie » struck), corrected in his ink: his comments to the
%     redactor of an exposé, par. 4 « Homomorphisme de Gysin » (4.1)-(4.4),
%     the list a)-j) of properties the formulaire must contain (projection
%     formula, trace of XVII, transitivity, Künneth, base change, change of
%     ring, independence of S, exchange formula, adjunction, linearity for
%     Jouanolou), a plan 4.1-4.3, the virtual dimension d(f), degree N with
%     Tr_f(1_X) = N 1_Y; then par. 5 « Classe fondamentale locale »
%     (4.4)-(4.9), purity, γ_{X/Y}, properties a)-h) and a « Remords » i)
%     with a triangle diagram. Item letters were re-ordered in ink; the
%     transcription gives the ink letters and notes the moves.
%   Page 10 — pencil, on an otherwise blank verso: the computation
%     (Σ d_i f_{i*} u_i^*) f^* = Σ d_i n_i id that page 11's last claim rests on.
%   Pages 15-20 — six slips of corrections to an exposé XVIII (1-1, 1.4, 1-2,
%     1-4, (1.4)//2, 1 2): Aut'(x) in pch(K), R^1 p_* p^* G, Γ(S,R^1p_*G),
%     a Mac Lane reference. Some slips (15, 17) are faint and read in part.
%     A second hand ticks « ok » in their margins; recorded in notes.
%
% Pages 1 (blank folder cover), 3 and 5 (covers, which become the section
% headings) are absent.
\documentclass[11pt,a4paper]{article}
% Le préambule commun à toutes les transcriptions du fonds Grothendieck.
%
% Il tient en une idée : **l'incertitude se compose, elle ne se résout pas.**
% Une transcription qui lisse un mot douteux a détruit la seule chose pour
% laquelle on la faisait — un lecteur doit pouvoir savoir, à chaque endroit,
% ce qui était sur le feuillet et ce qui a été ajouté. Les six macros qui
% suivent sont donc l'appareil critique, et non de la décoration.
%
% Les mêmes commandes sont reconnues par `scripts/render.mjs`, qui produit la
% vue de lecture du site. Une macro ajoutée ici doit l'être aussi là-bas,
% faute de quoi le rendu échouera bruyamment — ce qui est voulu : un rendu
% silencieusement faux est pire qu'un rendu absent.

\ProvidesPackage{grothendieck}[2026/08/08 Transcriptions du fonds Grothendieck]

\usepackage{amsmath,amssymb,amsthm}
\usepackage{mathrsfs}
\usepackage{stmaryrd}
% Les diagrammes commutatifs sont partout dans ce fonds : c'est la langue
% même de la géométrie algébrique de Grothendieck, et une transcription qui
% les rendrait en prose ne transcrirait rien.
\usepackage{tikz-cd}
% Français par défaut — c'est la langue des feuillets — mais l'anglais est
% chargé aussi : la traduction et le résumé sont des documents anglais, et la
% typographie française (espaces avant les deux-points) y serait une faute.
% Ces documents basculent par \selectlanguage{english} en tête de corps.
\usepackage[english,french]{babel}
\usepackage{fontspec}
\usepackage{geometry}
\geometry{a4paper,margin=25mm,marginparwidth=28mm}
\usepackage{marginnote}
\usepackage{xcolor}
% ulem, not soul. `soul`'s \st and \ul are built on a character-by-character
% scan that XeLaTeX's font machinery breaks: the document fails with an
% undefined control sequence rather than degrading. ulem does the same two jobs
% and is Unicode-engine safe. `normalem` stops it hijacking \emph, which the
% transcriptions use for Grothendieck's own emphasis.
\usepackage[normalem]{ulem}
\usepackage{hyperref}
\hypersetup{colorlinks=true,linkcolor=black,urlcolor=black,pdfborder={0 0 0}}

% --- Métadonnées du lot -----------------------------------------------------
% Reprises telles quelles par le script de rendu, qui les affiche en tête de
% la vue de lecture. Elles fixent ce que le fichier prétend couvrir : sans
% elles, une transcription flotte sans point d'attache dans un fonds de seize
% mille feuillets.

\newcommand{\folder}[1]{\gdef\grothFolder{#1}}
\newcommand{\batch}[1]{\gdef\grothBatch{#1}}
\newcommand{\leaves}[2]{\gdef\grothFirst{#1}\gdef\grothLast{#2}}
\newcommand{\foldertitle}[1]{\gdef\grothFoldertitle{#1}}
\gdef\grothFolder{?}\gdef\grothBatch{?}\gdef\grothFirst{?}\gdef\grothLast{?}\gdef\grothFoldertitle{}

% --- Le résumé --------------------------------------------------------------
%
% Il ouvre la lecture modernisée, sous le titre « Résumé », et il est composé
% en petit. Ce n'est pas un ornement : c'est le seul passage du document qui
% ne soit pas des mathématiques, et un lecteur doit pouvoir le reconnaître
% avant de l'avoir lu — puis le sauter s'il connaît déjà le sujet, ou s'y
% arrêter s'il ne le connaît pas. Le filet qui le suit marque la reprise du
% corps.

\newenvironment{resume}
  {\small\setlength{\parskip}{0.45em}}
  {\par\medskip\hrule\medskip}

% --- Mots-clés ---------------------------------------------------------------
%
% En anglais, en clôture du résumé de la lecture modernisée : le vocabulaire
% moderne sous lequel chercher ce que les feuillets construisent. Le manifeste
% du site (`npm run manifest`) les extrait pour étiqueter la cote entière —
% cette ligne est la source unique des tags, il n'y a pas de fichier à côté.
\newcommand{\keywords}[1]{\par\smallskip\noindent
  {\footnotesize\textsc{Keywords} --- \textit{#1}}\par}

% --- L'appareil critique ----------------------------------------------------

% Le numéro de feuillet, en marge. C'est l'ancre qui permet de régler un
% désaccord sur une formule en regardant *un* feuillet plutôt qu'une chemise
% de six cents. Le site s'en sert aussi pour tourner les pages du fac-similé
% quand on fait défiler la transcription.
\newcommand{\leaf}[1]{%
  \marginnote{\footnotesize\color{gray!70}\textsf{#1}}%
  \label{leaf:#1}\ignorespaces}

% Illisible : on ne devine pas. Un mot inventé qui se lit comme les autres est
% le pire résultat possible.
\newcommand{\ill}{\textcolor{red!60!black}{[\,\ldots\,]}}

% Lecture douteuse : le mot est donné, et signalé comme incertain.
\newcommand{\uncertain}[1]{\uline{#1}}

% Ajout de l'éditeur — une lettre manquante, un mot sous-entendu.
\newcommand{\add}[1]{\textcolor{blue!50!black}{[#1]}}

% Biffé par Grothendieck lui-même. Ce qu'il a rayé fait partie du document :
% une rature dit souvent où la pensée a changé de direction.
\newcommand{\struck}[1]{\sout{#1}}

% Note du transcripteur, sur l'état matériel du feuillet.
\newcommand{\note}[1]{\par\smallskip{\small\color{gray!60!black}\textit{[#1]}}\par\smallskip}

% Note marginale de Grothendieck, distincte du corps du texte.
\newcommand{\marginal}[1]{\par\smallskip\noindent\hspace{1em}%
  {\small\color{gray!60!black}#1}\par\smallskip}

% --- Titre ------------------------------------------------------------------

\AtBeginDocument{%
  \begin{center}
    {\large\bfseries Fonds Alexandre Grothendieck --- cote n\textsuperscript{o}~\grothFolder}\\[2pt]
    {\itshape\grothFoldertitle}\\[4pt]
    {\small feuillets \grothFirst--\grothLast\ (lot \grothBatch)}\\[2pt]
    {\footnotesize\color{gray!60!black}Transcription établie d'après le fac-similé numérisé
     par l'Université de Montpellier.\\
     Les lectures incertaines sont soulignées, les illisibles notées [\,\ldots\,],
     les ajouts entre crochets.}
  \end{center}
  \vspace{1em}}

\endinput


\folder{32}
\batch{1}
\pages{1}{20}
\dating{[vers 1963-1977]}
\watermark{Édition de démonstration}
\foldertitle{Corrections (références, rajouts) [dont SGA 4] : tapuscrit (s.d.), notes manuscrites (s.d.)}

\begin{document}

\section*{SGA 4 IX : lignes à corriger (page 2)}

\page{2}
\note{Feuillet à l'encre, seul en haut à gauche d'une page par ailleurs
vierge : un relevé de pages et de lignes de l'exposé IX de SGA 4, sans le
texte des corrections. « $\ell.\,-4$ » compte depuis le bas de la page.}

SGA 4 IX \note{« IX » est encadré.}

\begin{itemize}
\item[p.\ 2] $\ell.\,-4$
\item[3] $\ell.\,1$, $\ell.\,8$, $\ell.\,11$
\item[] \struck{\ill{}}
\item[60] $\ell.\,9$, $\ell.\,12$
\item[61] $\ell.\,\struck{1}\,4$
\item[\uncertain{63}] $\ell.\,12$
\item[80] $\ell.\,5$, $\ell.\,6$, $\ell.\,9$
\item[91] $\ell.\,-5$
\end{itemize}

\section*{\struck{Notes et lettres} XVII \uncertain{5.5} (page 4)}

\note{Le titre est celui de la chemise orange qui précède (page 3), de sa
main : « \struck{Notes et lettres} » biffé, puis « XVII » encadré, et
au crayon un numéro qu'on lit « 5.5 » sans certitude. La page 4, au crayon,
porte sur les foncteurs dérivés $LT$ d'un foncteur non additif $T$
(puissances divisées, extérieures, symétriques) via Dold--Puppe.}

\page{4}
$\Gamma^n$, $\Lambda^n$, $\mathrm{Sym}^n$, $\overset{\sim}{\otimes}$
\note{le dernier symbole est un $\otimes$ surmonté d'un signe qu'on lit
comme un tilde, sans certitude}

$D^{\leqslant 0}$ : commutation aux $\varinjlim$ [\uncertain{par passage
du complexe plat} \uncertain{donné par} Dold--Puppe\,; \ill{} \uncertain{il
vaut fibre par fibre}, et \ill{} O.K.\ \ill{} aux $\varinjlim$]

$D^{b,\,\mathrm{tor}\leqslant 0} \longrightarrow D^{\leqslant 0}$ :
\uncertain{mes ingrédients}

Si $\operatorname{tordim} K < -d$, alors $H^i(LT(K)) = 0$ si $i < d$.
\note{cette ligne est séparée de la suivante par un trait ; le $d$ de
« $-d$ » est surchargé}

Si degré $T \leqslant n$, alors si $H^i(K) = 0$ si $i > d$,
\[
H^i(LT(K)) = 0 \ \text{si}\ i \geqslant nd .
\]

Si $\operatorname{tordim} K < -d$, $\Rightarrow$ $\operatorname{tordim}
LT(K) < -d$ si $T$ transforme plats en plats.


Si $K$ dans intervalle $[p, q]$ \struck{\ill{}}, degré $T
\leqslant n$, $T(0) = 0$\,:
\[
N T N^{-1} K \quad
\begin{cases}
\text{si } 0 \leqslant p \leqslant q, & K \text{ dans } [p, nq]\\
\text{si } p \leqslant q \leqslant 0, & K \text{ dans } [np, q]
\end{cases}
\quad \text{comme complexes.}
\]
\note{au second membre, c'est $NTN^{-1}K$, et non $K$, qui est dans
l'intervalle indiqué ; la page écrit « $K$ ».}

Si $T$ conserve \uncertain{propriété} $P$ stable par facteurs directs,
\ill{}, alors si les comp.\ de $K$ sont $P$, idem $NTN^{-1}K$ …


\uncertain{Ajouter remarque formulaire.}

\section*{Notes \struck{Bibliographie} (pages 6 à 14)}

\note{Le titre est celui de la chemise qui précède (page 5), de sa main :
« Notes », entouré, et « Bibliographie » biffé. Les pages 6 à 14 sont un
tapuscrit : les paragraphes 4 (Homomorphisme de Gysin) et 5 (Classe
fondamentale locale) d'un texte sur la dualité et les classes
fondamentales, où l'auteur s'adresse au rédacteur (« je suggère », « Je
pense que », « Remords ») ; il cite SGA 2, SGA 5 IV, « XVI 3.7 », et
Jouanolou. Les corrections à l'encre sont de sa main : on les donne comme
ajouts et biffures ; les fautes de frappe corrigées sur la machine (xxx) ne sont pas
relevées, sauf quand le mot biffé se lit et compte.}

\page{6}
\subsection*{4. Homomorphisme de Gysin.}

Considérons un diagramme commutatif

\begin{tikzcd}
X \arrow[rr, "f"] \arrow[dr, "p"'] & & Y \arrow[dl, "q"] \\
& S &
\end{tikzcd}

(4.1), avec $p$ et $q$ compactifiables, \add{$p$ plat de prés.\
finie, $q$ lisse,} $S$ étant muni d'un Anneau de torsion $A$ premier aux
car.\ résiduelles. Pour un complexe de $A$-modules $F$ sur $S$, on définit
alors un homomorphisme canonique
\[
\text{(4.2)}\qquad \add{\mathrm{Tr}_f \text{ ou }} f_{*}\colon\;
Rf_{!}(F_X) \longrightarrow F_Y(d'-d)[2(d'-d)] ,
\]
en supposant $p$, $q$ de dimensions relatives $d$, $d'$ en tout point (il
suffit que $p$ soit de dim rel $\leqslant d$ en tout point). Si $p$ est
également lisse et de dim.\ rel.\ exactement $d$ partout, alors
l'homomorphisme précédent s'identifie à l'homomorphisme d'adjonction pour
$Rf_{!}$, $Rf^{!}$ relatif au deuxième membre. En tous cas, l'homomorphisme
en question s'appelle \emph{homomorphisme de Gysin} (relatif à $f$, $F$), ou
\emph{homomorphisme trace} (parfois noté $\mathrm{Tr}_f$). On appelle encore
par le même nom, et on note de même, tout homomorphisme déduit de (4.2) par
application d'un foncteur quelconque (le \struck{cas échéant} \add{plus
souvent,} en transformant encore la flèche obtenue en appliquant des
isomorphismes canoniques, tels des isomorphismes de transitivité …,
utilisés communément comme morphismes d'identification …). Par exemple on a
les homomorphismes
\[
\text{(4.2.1)}\qquad H^i(f_{*}) \text{ ou } f_{*}\colon\;
R^i f_{!}(F_X) \longrightarrow
\underline{H}^{i\supplied{+2(d'-d)}}(F_Y)(d'-d)\,\struck{[2(d'-d)]} ,
\]
de faisceaux sur $Y$, déduit\supplied{s} de (4.2.1) en appliquant
$\underline{H}^i$, ou
\note{la page renvoie à « (4.2.1) » dans la définition même de (4.2.1) ; il
s'agit sans doute de (4.2)}
\[
\text{(4.2.2)}\qquad Rq_{!}(f_{*}) \text{ ou } f_{*}\colon\;
Rp_{!}(F_X) \longrightarrow Rq_{!}(F_Y)(d'-d)[2(d'-d)] ,
\]
déduit de (4.2) en appliquant $Rq_{!}$, ou plus généralement, pour tout
morphisme compactifiable
\[
g\colon Y \longrightarrow Z
\]

\page{7}
\[
\text{(4.2.3)}\qquad R(gf)_{!}(F_X) \longrightarrow
Rg_{!}(F_Y)(d'-d)[2(d'-d)] ,
\]
déduit de (4.2) en appliquant $Rg_{!}$. De même, appliquant $Rg_{*}$, on
trouve
\[
\text{(4.2.4)}\qquad Rg_{*}Rf_{!}(F_X) \overset{\mathrm{dfn}}{=}
R(gf)_{c(f)}(F_X) \longrightarrow Rg_{*}(F_Y)(d'-d)[2(d'-d)]
\]
(où la \struck{première} \add{source} est « l'image directe à supports
propres sur $Y$ »). \struck{Prenant en particulier} Pour définir cette
dernière flèche, quitte à faire rentrer le $(d'-d)$ dans la parenthèse,
\add{il} suffisait de disposer d'un morphisme de \struck{sites} topos
$Y_{\text{ét}} \to Z$, et prenant alors pour $Z$ le topos final, on trouve
l'exemple important d'homomorphisme de Gysin\,:
\[
\text{(4.2.5)}\qquad R\Gamma_X(Rf_{!}(F_X)) \overset{\mathrm{dfn}}{=}
R\Gamma_{X,\,c(f)}(F_X) \longrightarrow R\Gamma_Y(F_Y(d'-d))[2(d'-d)] .
\]
Dans le cas où $f$ est \emph{propre}, (4.2) devient
\[
\text{(4.2.6)}\qquad Rf_{*}(F_X) \longrightarrow F_Y(d'-d)[2(d'-d)] ,
\]
et (4.2.4) et (4.2.5) s'écrivent respectivement
\[
\begin{aligned}
&\text{(4.2.7)}\qquad R(gf)_{*}(F_X) \longrightarrow Rg_{*}(F_Y)(d'-d)[2(d'-d)] ,\\
&\text{(4.2.8)}\qquad R\Gamma_X(F_X) \longrightarrow R\Gamma_Y(F_Y(d'-d))[2(d'-d)] .
\end{aligned}
\]
\note{dans (4.2.4) et (4.2.5), l'indice « $!$ » tapé d'abord est biffé et
remplacé par « $c(f)$ » ; au premier membre de (4.2.5) le tapuscrit porte
$R\Gamma_X$ là où, $Rf_{!}(F_X)$ vivant sur $Y$, on attendrait
$R\Gamma_Y$.}

\struck{Dans (4.2.5), (4.2.8), il y aurait lieu de remplacer $F_X$ par un
$F_X(i)$} On a également à définir des homomorphismes de Gysin pour des
coéfficients sur $Y$ plus généraux que ceux de la forme $F_Y$ et ceux qu'on
en déduit par twist et shift. De façon générale, soit $G$ un complexe de
$A$-Modules (NB je ne précise pas ici \add{ni ailleurs} les hypothèses de
degré) sur $Y$, et appliquons à (4.2) le foncteur
$\overset{L}{\otimes}\, G$, on trouve un homomorphisme canonique, qui est
d'ailleurs le même que si on partait de $F' = A_S$ sur $S$, et qu'on
tensorise l'homomorphisme de Gysin correspendant par
$F_Y \otimes G$. Donc, autant oublier $F$ et dire que pour tout $G$ sur
$Y$, on a un homomorphisme canonique (dit \add{\uncertain{encore}} de Gysin)
\[
\text{(4.3)}\qquad \boxed{\,Rf_{!}(f^{*}(G)) \longrightarrow
G(d'-d)[2(d'-d)]\,} ,
\]
ou, au choix
\marginal{Question\,: a-t-on besoin d'un $A$ sur $S$, ne suffit-il d'avoir
un Anneau $A_Y$ sur $Y$ de torsion, premier à \uncertain{$\ell$}\,?}

\page{8}
\[
\text{(4.3 bis)}\qquad Rf_{!}(f^{*}G(d-d')) \longrightarrow G[2(d'-d)] .
\]
Arrivé à ce point, je constate que c'est par cette flèche \add{(4.3)} (plus
raisonnable que (4.2)) qu'il aurait fallu commencer, en ajoutant encore que
si $p$ est également lisse et est de dim.\ rel.\ $d$ partout, et si $G$ est
localement (sur $Y_{\text{ét}}$) isomorphe à l'image inverse d'un complexe
sur $S$, alors (4.3) s'identifie à l'homomorphisme d'adjonction qu'on pense.
Il y aurait lieu alors de donner les exemples (4.2.1) à (4.2.8) en partant
de (4.3) plutôt que de (4.2). On peut laisser tomber (4.2.1), et se borner
à dire que les homomorphismes déduits des homomorphismes envisagés par
application des foncteurs $H^i$ sont également affublés du nom « Gysin » ou
« trace ».

Une perplexité\,: dans le contexte actuel, la notation $f_{*}$, qui dénote
à la fois un foncteur (image directe) et un homomorphisme de foncteurs
(Gysin) ne semble pas bien fameux. Je propose de donner la préférence ici à
la notation $\mathrm{Tr}_f$, en le disant au moment ou on signale également
la notation $f_{*}$.

Etant donné le nombre d'exemples parmi les variantes à Gysin, il est clair
qu'il ne sera pas possible de donner un formulaire qui donne\supplied{rait,}
\struck{immédiatement} sans déductions du tout, tous les cas qu'on
rencontrera. Il faudra se borner à donner un formulaire directement en
termes des homomorphismes de Gysin fondamentaux sous la forme (4.3), et se
borner à donner comme exemple la forme que prend ce formulaire « après
application d'un foncteur $Rg_{*}$ » dans le cas $f$ propre, disons.

Voici des points qui doivent certainement figurer au formulaire\,:

\marginal{descend}
\add{d)} Compatibilité avec la formule de projection\,: Si on a un
complexe de modules $G'$ sur $Y$ et si on tensorise (4.3) par $G'$, en
utilisant la formule de projection de XVII …, on trouve Gysin pour
$G \overset{L}{\otimes} G'$. Ceci \struck{montre} \add{entraîne} en
particulier que, via la formule de projection,
\note{les lettres des articles de cette liste sont récrites à l'encre
sur celles du tapuscrit ; on donne les lettres à l'encre. La formule de
projection, d'abord en tête, devient d) ; le mot « descend » dans la marge
et un trait tiré sous l'article c) à la page suivante marquent où elle doit
aller.}

\page{9}
(4.3) est connu pour tout $G$ par la seule connaissance de
\[
\text{(4.4)}\qquad Rf_{!}(A_X(d-d')[2(d-d')]) \longrightarrow A_Y .
\]
Ceci devrait bien donner la formule de projection au niveau des classes de
cohomologie (qu'il faut bien entendu expliciter, soit ici, soit \add{(mieux)}
à la fin du formulaire, en reprenant ce formulaire dans le cas particulier
recommandé\supplied{,} parmi les variantes possibles).

\add{a)} \emph{Si} $d = d'$ \emph{et si} $f$ \emph{est quasi-fini} (et
alors automatiquement plat, supposant que la dimension relative de $p$ est
partout $d$), alors on retrouve l'homomorphisme trace de XVII …

\add{b)} Si $f$ \emph{est lisse}, on retrouve l'homomorphisme trace du
par.\,2.

\add{c)} \emph{Transitivité} pour un composé $gf$ ($g\colon Y \to Z$, avec
$Z$ lisse et compactifiable sur $S$).
\note{sous cet article, un trait à l'encre venu de la marge : c'est ici que
doit venir l'article d) de la page 8.}

\add{f)} Compatibilité avec \emph{Künneth}, quand on regarde
$g_1 \times_S g_2\colon X_1 \times_S X_2 \to Y_1 \times_S Y_2$.
Compatibilité avec \emph{changement de base} $S' \to S$.

\add{e)} Compatibilité avec \emph{changement d'Anneau} $A \to A'$ (qui
montre qu'il suffit de connaitre (4.4) pour les
$A = (\mathbb{Z}/n\mathbb{Z})_S$ pour l'avoir pour tout $A$).

g) \emph{Indépendance par rapport à} $S$ (si on se donne $S \to T$ lisse,
de sorte que l'on peut regarder $X$ et $Y$ comme étant sur $T$ …).

h) \emph{Formule d'échange}, relatif au cas où on a un diagramme cartésien

\begin{tikzcd}
X \arrow[d, "f"'] & X' \arrow[l, "h"'] \arrow[d, "f'"] \\
Y & Y' \arrow[l, "g"]
\end{tikzcd}

de schémas sur $S$\,: appliquant $g^{*}$ à $\mathrm{Tr}_f$, et
appliquant l'isomorphisme de changement de base pour $Rf_{!}$ et $g^{*}$,
on doit trouver $\mathrm{Tr}_{f'}$. On supposera bien sûr $Y$ et $Y'$
lisses sur $S$, $X$ plat sur $S$, et pour que $X'$ soit plat sur $S$, il
faut évidemment une hypothèse genre Tor-indépendance de $X$ et $Y'$ sur
$Y$. Cela doit marcher en tous cas si $X$ est également lisse sur $S$, et
si $f$ et $g$ sont \emph{transversaux} (EGA IV 17 …).
\note{le carré est entouré au crayon d'un trait fermé.}

\page{10}
\note{Calcul au crayon dans la moitié inférieure d'une page par ailleurs
vierge, sans lien écrit avec le tapuscrit : le degré d'un morphisme fini
$f$ lu sur la cohomologie, $f_{*}f^{*} = N$, à partir des points $x_i$ de
la fibre et de leurs adhérences $Z_i$.}

$Z_i = \overline{x_i}$ \qquad $X_y \ni x_i$ \qquad $y$ \qquad
$\deg(x_i/y) = n_i$
\[
\textstyle\sum d_i n_i = N
\]

\begin{tikzcd}
X \arrow[r, "f"] & Y \\
Z_i \arrow[u, "u_i"] \arrow[ur, "f_i"'] &
\end{tikzcd}

\begin{tikzcd}
H^r(X) \arrow[dr, "u_i^{*}"'] & & H^r(Y) \arrow[ll, "f^{*}"'] \\
& H^r(Z_i) \arrow[ur, "f_{i*}"'] &
\end{tikzcd}

\[
\Bigl(\sum d_i\, f_{i*} u_i^{*}\Bigr) f^{*} = \sum d_i\, f_{i*} u_i^{*} f^{*}
= \sum d_i\, (f_{i*} f_i^{*}) = \sum d_i\, n_i\, \mathrm{id}
\]
\note{le libellé de la flèche $H^r(Y) \to H^r(X)$ est surchargé ; on le lit
$f^{*}$. Dans la dernière égalité, la page écrit « $\sum d_i (f_{i*}f_i^{*})
n_i\, \mathrm{id}$ », le facteur $n_i$ venant après la parenthèse
entourée.}

\page{11}
i) \emph{Formule de \struck{adjonction} transposition ou d'adjonction},
disant que Gysin sur la cohomologie à supports propres est transposé de
l'homomorphisme image inverse. J'ai un sentiment que cela résulte
tautologiquement de la définition, \struck{et n'a guère de sens que dans}

j) \emph{Formules de linéarité} (notamment pour $f$ une immersion de
schémas relatifs lisses)\add{,} nécessaire pour Jouanolou\,; regarder aussi
la formule de linéarité pour variétés éclatées déduites de cette
situation. Si possible, donner formulation commune\,?

Comme plan possible du par.\,4, je vois
\begin{itemize}
\item[4.1.] Définition et notations.
\item[4.2.] Formulaire général.
\item[4.3.] Formulaire spécial, pour classes de cohomologie.
\end{itemize}

Du point de vue notations, je suggère, pour ne pas trainer indéfiniment
des $d-d'$ dans les notations, de travailler systématiquement avec la
\emph{dimension virtuelle} $d(f) = d-d'$ de $f$\,; c'est elle qui mérite une
lettre simple, telle que $d$.

Dans 4.3, inclure le cas $f$ propre, $d = d'$\,; dire que dans ce
\struck{5. Classes fondamentales locales} cas on a
\[
\mathrm{Tr}_f(1_X) = N\, 1_Y
\]
où $N$ est une fonction entière localement constante sur $Y$, appelée
\emph{degré} de $f$, et qu'on peut caractériser comme étant égale, sur le
plus grand ouvert de $Y$ sur lequel $f$ soit fini localement libre, comme
égale au degré habituel. Conclure que $\operatorname{Ker}(f^{*}\colon
H^{*}(Y, G) \to H^{*}(X, G_X))$ est annulé par $N$, donc \add{$f^{*}$} est
injectif (et même inversible à gauche) si $N = 1$. [Si $S = \operatorname{Spec}
k$ et si $Y$ est irréductible, de point générique \add{$\eta$},
\add{($d$ et $d'$ quelconques)} et si $N$ est le pgcd des degrés des
$0$-cycles de $X_{\eta}$ \add{$\neq \varnothing$}, on peut de même trouver
un homomorphisme\supplied{,} fonctoriel en $G$\add{,} $H^{*}(X, G_X) \to
H^{*}(Y, G)$ dont le composé avec l'homomorphisme image inverse est
$N.\mathrm{id}$.]
\note{les crochets, et un long trait de crayon qui court dans la marge
gauche depuis « $\mathrm{Tr}_f(1_X) = N\,1_Y$ », sont de sa main. Le
calcul de la page 10 est celui qui justifie ce dernier énoncé.}

\page{12}
\subsection*{5. Classe fondamentale locale.}

Reprenons l'homomorphisme de Gysin général (4.3), et notons que par la
formule d'adjonction, sa connaissance équivaut à celle d'un homomorphisme
\[
\text{(4.4)}\qquad f^{*}(G)(d-d')[2(d-d')] \longrightarrow Rf^{!}G ,
\]
\marginal{autre nom\,?}
[On l'appelle l'homomorphisme \emph{classe fondamentale locale} (?)] Il est
fonctoriel en $G$, et, comme nous avons vu au par.\,4, il suffit de le
connaitre pour $G = A_Y$ pour le connaitre pour tout $G$. Dans ce cas, on
peut l'interpréter comme un homomorphisme
\[
\text{(4.5)}\qquad \varphi_f\colon A_X \longrightarrow
Rf^{!}(A_Y)(d'-d)\ 2(d'-d)\add{)} ,\quad
\add{\text{i.e. } \varphi_f \in H^{2(d-d')}(Y, Rf^{!}(A_Y)(d-d'))}
\]
et il définit par suite une section, notée de même
\[
\text{(4.6)}\qquad
\begin{aligned}
&\struck{\varphi_f \in H^0(Y, \underline{H}^0 Rf^{!}(A_Y)(d'-d)\ 2(d'-d))}\\
&\add{\varphi_f} \in H^0(Y, \underline{R}^{2(d'-d)} f^{!}(A_Y))(d'-d)) ,
\end{aligned}
\]
qu'on appelera également (et peut-être à meilleur escient) la \emph{classe
fondamentale locale pour $f$}. On notera que pour la connaitre pour tout
$A$, il suffit de la connaitre pour les $\mathbb{Z}/n\mathbb{Z}$ (du moins
si on se borne aux $A$ qui sont annulés par un même entier $n \geqslant 1$
premier aux car.\ résiduelles de $S$)\,: \struck{La connaissance de (4.6)
permet de reconstituer} plus précisément, la formation des classes
$\varphi_f$ de (4.5) et de (4.6) est \emph{fonctoriel} en l'Anneau $A_S$
(voire, en $A_Y$).
\note{(4.4) porte ici le numéro déjà pris au par.\,4, et la suite du
par.\,5 le garde. Dans (4.5) et (4.6) le tapuscrit n'a pas les crochets
du décalage ; la parenthèse fermante après « $2(d'-d)$ » de (4.5) est
ajoutée à l'encre, et l'ajout « i.e. … » qui suit, surchargeant une ligne
biffée, porte $2(d-d')$ et $(d-d')$ avec les signes inverses de ceux de
la ligne tapée. Au premier membre de l'ajout de (4.6), le symbole est le
$\varphi$ de sa main ; dans le tapuscrit c'est un $\emptyset$ barré.}

Lorsqu'on dispose d'une « formule de pureté »
\[
\text{(4.7)}\qquad R^i f^{!}(A_Y) = 0 \quad \text{pour } i \neq 2(d-d') ,
\]
la connaissance de (4.6) équivaut même à celle de (4.5), ce dernier
$\varphi_f$ méritant le nom de \emph{classe fondamentale globale} de $X$
sur $Y$.
\note{dans (4.7), le signe entre $i$ et $2(d-d')$ est surchargé à l'encre ;
on lit $\neq$, sans certitude.}

Nous allons nous intéresser par la suite au cas où $f\colon X \to Y$ est
une \emph{immersion fermée}. On sait alors que les formules de pureté (4.7)

\page{13}
sont vérifiées. (Lorsque $X$ est également lisse, c'est contenu dans
XVI 3.7\,; le cas général --- sans même supposer $Y$ plat sur $S$,
d'ailleurs --- s'en déduit par dévissage facile, et il est traité d'ailleurs
dans SGA 2 XIV 1, conjonction de 1.10 et de 1.18.) La classe $\varphi_f$ de
(4.5) peut alors s'interpréter aussi comme un élément
\[
\text{(4.8)}\qquad \add{\gamma_{X/Y}} \in H^{2(d-d')}_{\add{X}}(Y, A_Y(d-d')) .
\]
Cette classe s'appelle la \emph{classe de cohomologie à supports dans $X$
de $X$ dans $Y$}, et son image
\[
\text{(4.9)}\qquad \add{\gamma_{X/Y}} \in H^{2(d-d')}(Y, A_Y(d-d'))
\]
\emph{la classe de cohomologie} (sous-entendue\,: à supports quelconques)
\emph{de $X$ dans $Y$}. Le formalisme des classes de cohomologie
\add{(4.8), (4.9)} ainsi \struck{associées} \add{obtenues} (une fois la
définition étendue par linéarité à des cycles algébriques relatifs sur $Y$)
sera étudié systématiquement dans SGA 5 IV. Notre propos ici est de donner
les propriétés les plus élémentaires de la classe locale (4.6), en attachant
notre attention surtout au cas où $X$ est également lisse sur $S$.
\note{dans (4.8), le tapuscrit portait $\varphi_f$ et $H^{2(d-d')}(Y, \ldots)$ ;
$\gamma_{X/Y}$ et l'indice $X$ sont de sa main, écrits par-dessus. Dans
(4.9), $\gamma_{X/Y}$ est de même à l'encre.}

Voici une liste de propriétés à expliciter qui me semble raisonnable (où on
pose \struck{$p$} \add{\uncertain{$c$}} $= d-d'$)\,:
\note{la lettre $p$ du tapuscrit est récrite à l'encre ici et aux articles b)
et c) par un signe bouclé qu'on lit $c$ (la codimension), sans certitude.}

a) Si $X$ est lisse sur $S$, et si $G$ est localement \add{(ét)} sur $Y$
isomorphe à l'image inverse d'un complexe sur $S$, \add{p.\,ex.\
$G = (A)_Y$,} alors (4.4) est un isomorphisme. (C'est essentiellement
contenu dans XVI 3.7).

b) Si \add{\uncertain{$c$}} $= 0$ i.e.\ $f$ un isomorphisme, (4.4) est
l'identité i.e.\ $\varphi_f = 1$ dans (4.5) ou (4.6).

c) Si \add{\uncertain{$c$}} $= 1$, de sorte que $X$ est un diviseur relatif
sur $Y/S$, alors $\varphi_f$ de (4.5) i.e.\ (4.8), et a fortiori de (4.6)
ou (4.9), est donné par la théorie de Kummer comme d'habitude.
\note{la lettre « c) » est entourée au crayon.}

d) Compatibilité de $\varphi_f$ avec \emph{transitivité} et
\emph{intersections} (ces deux formules sont essentiellement identiques),
dans le cas \emph{lisse}

\page{14}
\emph{et transversal}.

e) Compatibilité avec \emph{changement de base} $S' \to S$, avec
\emph{produits cartésiens} $Y_1 \times_S Y_2$ sur $S$.

g) Indépendance par rapport à $S$. \add{($S \to T$)}

f) Compatibilité avec changement d'anneau.
\note{la lettre g) surcharge une autre lettre tapée ; une flèche au crayon
dans la marge fait passer l'article g) après f).}

h) Relation avec Gysin\,: Les classes (4.8), (4.9) s'interprètent comme
l'image de $1_X$ par Gysin. Le transformé de (4.4) par $Rf_{!}$, suivi de
l'homomorphisme d'adjonction, est l'homomorphisme de Gysin.
\note{un trait de crayon part d'ici, descend le long de la marge gauche et
aboutit à l'article « i) » du « Remords » plus bas : c'est là qu'il doit
venir, à la suite de h).}

Je pense que, compte tenu du formulaire du par.\,4, le formulaire proposé
ici pour les classes fondamentales $\pm$ locales est pratiquement trivial à
partir des définitions. On peut, pour finir, récupérer un énoncé qui donne
la \emph{caractérisation} des classes fondamentales locales dans le cas d'un
couple lisse, par les propriétés suivantes 1) Nature locale (pour Zar
suffit) 2) cas $p = 0$ 3) cas $p = 1$, quand $X$ est défini par une
équation 4) propriété de cup-produit d) (où il suffit de prendre un des
deux facteurs de codimension 1). Signaler \add{en remarque} que cet énoncé
\add{(existence et unicité d'une théorie des cl.\ fond.\ loc.)} peut se
démontrer sans difficulté indépendamment de la théorie de dualité et de
celle de l'homomorphisme trace, et avait été donnée dans le séminaire oral
avant la dualité.

Remords\,: il serait prudent d'adjoindre au formulaire un

i) On a commutativité dans le diagramme suivant

\begin{tikzcd}[column sep=large]
f^{*}(G)(\supplied{-}p)[\supplied{-}2p] \arrow[r] \arrow[dr, bend right=20] &
Rf^{!}(G) \\
& Rf^{!}(A_Y) \overset{L}{\otimes} f^{*}(G) \arrow[u]
\end{tikzcd}

où la flèche horizontale est (4.4), la flèche oblique est
déduite par $\overset{L}{\otimes}\, G$ de (4.5), et le flèche verticale est
l'homomorphisme canonique qu'on devine, analogue à celui de changement de
base pour les ${}^{!}$, et qui manque dans 3.1 (tout comme ce dernier). On
s'y intéressera sans doute surtout si $G$ est localement libre, donc
l'homomorphisme en question est un isomorphisme.
\note{les signes « $-$ » du premier terme du diagramme sont ajoutés à
l'encre ; ici la lettre $p$ du tapuscrit n'est pas récrite.}

\section*{Corrections à l'exposé XVIII (pages 15 à 20)}

\note{Six fiches de format irrégulier, à l'encre bleue ou au crayon,
chacune en tête d'un renvoi « XVIII » (souligné ou encadré) suivi d'un
numéro de paragraphe ; les renvois de pages et de lignes (« $\ell.\,9$ du
bas ») visent un tapuscrit de l'exposé XVIII qui n'est pas ici. Le
vocabulaire (images directes, catégories dérivées, $\mathrm{pch}(K)$,
$R^1p_{*}$) est celui des champs de Picard et de la cohomologie des fibrés
projectifs. Dans la marge gauche de chaque article, une autre main a écrit
au crayon noir « ok » ou un signe raturé ; on les donne en note.}

\page{15}
XVIII 1-1 \note{« XVIII » est encadré, ici et à la page 17.}

p.\ 2.

$\ell.$ 2 du haut \note{« ok » dans la marge}

\ill{} à la \ill{} \uncertain{commutant aux images directes}.
\struck{\ill{} ce qui \ill{} \uncertain{appliquer les catégories dérivées}.}
\note{les deux dernières lignes sont biffées d'un gribouillis au crayon.}


$\ell.$ 9 des bas \note{dans la marge, un signe raturé et un « ! »}

\uncertain{pourquoi supposer $X$ cohérente à présent}\,?


p.\ 3, $\ell.$ 7-8 du haut \note{« ok » dans la marge}

\uncertain{embrouillé} ---

\page{16}
1.4. \uncertain{14}

$\ell.$ 4 du bas

Peut-être par le \uncertain{faisceau} $\underline{\mathrm{Aut}}'(x)$ \,?

des automorphismes de $x$ qui sont \emph{dans} $\mathrm{pch}(K)$\,:
\note{en haut à droite, « 1.4. » et, souligné, un « 14 » qu'on lit sans
certitude ; le $x$ de $\mathrm{Aut}'(x)$ est surchargé.}

\page{17}
XVIII --- 1-2

$\ell.$ 7 : $K_0$ \ill{} \uncertain{mais} $G$
\note{« ok » dans la marge, ici et aux deux articles suivants}
\[
K(K_0, \varphi) = p^{*}(K_0) + p^{*}(\varphi)[\mathcal{O}(1)]
\]


lig.\ 1 du bas

\uncertain{peut-être mettre} $R^1 p_{*} p^{*} G$


1.2.2.1\,: \uncertain{j'ai eu du mal à comprendre d'abord lire le dernier}
\ill{}


\uncertain{on pourrait faire remarquer}
\[
\Gamma(S, R^1 p_{*} G) = \frac{H^1(P(E), G)}{H^1(S, G)}
\]
\uncertain{qu'il y a une section}
\note{cette fiche est au crayon, très pâle ; les lignes de prose ne sont lues
qu'en partie.}

\page{18}
1-4

$\ell.$ 9, 10, 11, 12 du bas ?
\note{« ok ? » dans la marge ; les quatre numéros de lignes sont réunis par
une accolade, et le point d'interrogation est à droite d'un trait vertical.}

\page{19}
XVIII --- (1.4)\,//\,2

$\ell.$ 1 du haut

on peut peut-être mettre une réf.

Mac-Lane

\page{20}
XVIII --- 1 2

p.\ 8

$\ell.$ 4 du haut \note{« ok » dans la marge}

réf 1.2.1 annulé

\end{document}
